\documentclass[11pt]{article}

%%METADATA
\title{GA1025 Macro Theory \\Solution to Problem Set 11}
\author{
Junbiao Chen\thanks{E-mail: jc14076@nyu.edu.}
}
\date{\today}


%%PACKAGES
\usepackage{mdframed} % For boxed environments
\usepackage{graphicx}
\usepackage{grffile}
\usepackage{tabularx}
\usepackage{setspace}
\usepackage{amsmath,amsthm,amssymb}
\usepackage[hyphens]{url}
\usepackage{natbib}
\usepackage[font=normalsize,labelfont=bf]{caption}
\usepackage[margin=1in]{geometry}
\usepackage{hyperref}
\hypersetup{colorlinks=true,urlcolor=blue,citecolor=blue}
\usepackage{stmaryrd}  %Package with \boxast command
\usepackage{enumerate}% http://ctan.org/pkg/enumerate %Supports lowercase Roman-letter enumeration
\usepackage{verbatim} %Package with \begin{comment} environment
%\usepackage{enumitem}
\usepackage{physics}
\usepackage{tikz}
\usepackage{listings}
\usepackage{upquote}
\usepackage{booktabs} %Package with \toprule and \bottomrule
\usepackage{etoc}     %Package with \localtableofcontents
\usepackage{placeins}    %Package that prevent repositioning the tables
\usepackage{multicol}
\usepackage{bm}
\usepackage{subfig}
\usepackage{csquotes}

\definecolor{dkgreen}{rgb}{0,0.6,0}
\definecolor{gray}{rgb}{0.5,0.5,0.5}
\definecolor{mauve}{rgb}{0.58,0,0.82}

\lstset{language=bash,
  frame=tb,
  aboveskip=3mm,
  belowskip=3mm,
  showstringspaces=false,
  columns=flexible,
  basicstyle={\small\ttfamily},
  numbers=none,
  numberstyle=\tiny\color{gray},
  keywordstyle=\color{blue},
  commentstyle=\color{dkgreen},
  stringstyle=\color{mauve},
  breaklines=true,
  breakatwhitespace=false,
  tabsize=3
}

\lstset{language=C,
  aboveskip=3mm,
  belowskip=3mm,
  showstringspaces=false,
  columns=flexible,
  basicstyle={\small\ttfamily},
  numbers=none,
  numberstyle=\tiny\color{gray},
  keywordstyle=\color{blue},
  commentstyle=\color{dkgreen},
  stringstyle=\color{mauve},
  breaklines=true,
  breakatwhitespace=false,
  tabsize=4
}

\definecolor{lightblue}{rgb}{0.68, 0.85, 0.9} %

%CUSTOM DEFINITIONS
\theoremstyle{definition}
\newtheorem{definition}{Definition}[section]
\newtheorem*{remark}{Remark}
\setcounter{secnumdepth}{3}
\usepackage{tikz}
\usetikzlibrary{arrows.meta}
\usetikzlibrary{automata, positioning, arrows, calc}

\tikzset{
	->,  % makes the edges directed
	>=stealth, % makes the arrow heads bold
	shorten >=2pt, shorten <=2pt, % shorten the arrow
	node distance=3cm, % specifies the minimum distance between two nodes. Change if n
	every state/.style={draw=blue!55,very thick,fill=blue!20}, % sets the properties for each ’state’ n
	initial text=$ $, % sets the text that appears on the start arrow
 }

%% PROPOSITION
% Define the Proposition environment
\newmdenv[
  innerleftmargin=10pt, 
  innerrightmargin=10pt,
  innertopmargin=10pt,
  innerbottommargin=10pt,
  linecolor=black, 
  linewidth=1pt,
  backgroundcolor=white, 
  roundcorner=5pt
]{propositionbox}

\newtheoremstyle{boldtitle} % Define a new theorem style
  {10pt} % Space above
  {10pt} % Space below
  {\itshape} % Body font
  {} % Indent amount
  {\bfseries} % Theorem head font
  {.} % Punctuation after theorem head
  { } % Space after theorem head
  {} % Theorem head spec

\theoremstyle{boldtitle} % Use the custom style
\newtheorem{proposition}{Proposition} % Define the proposition environment

% Redefine the proposition environment to use the box
\newenvironment{boxedproposition}[1][]
{\begin{propositionbox}\begin{proposition}[#1]}
{\end{proposition}\end{propositionbox}}

%%FORMATTING
\usepackage[bottom]{footmisc}
\onehalfspacing
\numberwithin{equation}{section}
\numberwithin{figure}{section}
\numberwithin{table}{section}
\bibliographystyle{../bib/aeanobold-oxford}


%main text
\begin{document}

\maketitle

\section{Exercise 10.8}
\subsection{(a) Answer.}
The Lagrangian of type 1 and 2 agents are the following, respectively.
\begin{align*}
    \mathcal{L}^1 & = \log c^t_t + \log c^t_{t+1} + \lambda ( y - c^t_t - \frac{p_{t+1} c^t_{t+1}}{p_t}) \\ 
    \mathcal{L}^2 & = \log c^t_t + \log c^t_{t+1} + \lambda ( \frac{p_{t+1} Y}{p_t} - c^t_t - \frac{p_{t+1} c^t_{t+1}}{p_t})
\end{align*}
The FOCs for Type 1 Youth are 
\begin{align*}
    [c^t_t] \quad & \frac{1}{c^t_t} = \lambda \\ 
    [c^t_{t+1}] \quad & \frac{1}{c^t_{t+1}} = \lambda \frac{p_{t+1}}{p_t}
\end{align*}
Combining the budget constraint, we have 
\[
{c_t^t}^{, \text{type 1}} = \frac{y}{2}
\]
Similarly, we have 
\[
{c_t^t}^{, \text{type 2}} = \frac{Y p_{t+1}}{2 p_t}
\]
As a result, the their saving functions are given by 
\[
{s_t^t}^{, \text{type 1}} = y - \frac{y}{2} = \frac{y}{2}
\]
\[
{s_t^t}^{, \text{type 2}} = 0 - \frac{Y p_{t+1}}{2 p_t} = - \frac{Y p_{t+1}}{2 p_t} = - \frac{Y}{2 R_t}
\]

\textbf{Since there is no trade between generations}, the consumption loan market clears within generations.
Thus, we have 
\[
\frac{y}{2} N_1 - \frac{Y p_{t+1}}{2 p_t} N_2 = 0 \Rightarrow y N_1 - \frac{Y N_2}{1 + r_t} = 0,
\]
which gives us the formula for the return as 
\[
1 + r(t) = \frac{Y N_2}{y N_1}.
\]


\subsection{(b) Answer.}
Assuming $p_t = q H_0$, then we have
\[
s(R_t) = \frac{H_0}{p_t} = \frac{H_0}{q H_0} = \frac{1}{q}
\]
Plugging in the aggregate saving function derived in (a), we have 
\begin{align*}
    \frac{y}{2} N_1 - \frac{Y p_{t+1}}{2 p_t} N_2 & = \frac{1}{q} \\ 
    \Rightarrow \frac{y}{2} N_1 - \frac{Y }{2 R_t } N_2 & = \frac{1}{q}
\end{align*}
Now imposing $R_t = 1$, we have 
\[
q = \frac{2  }{yN_1 - Y N_2},
\]
which is negative due to the assumption of $1 + r_t > 1 \Rightarrow y N_1 > Y N_2$.
It contradicts to the requirement of positive $q$.
Therefore, there is no quantity-theory- style equilibrium where fiat currency is valued.

\subsection{(c) Answer.}
Using the same equation as (b), but now we assume $1 + r_t < 1$,
then we have a quantity-theory- style equilibrium where fiat currency is valued and 
the $q$ is positive. 

\noindent \textbf{Remark:} Now, we examine the distributional effects under valued fiat money.
Since type 1 and type 2 agents are lenders and borrowers, respectively.
A change in the interest rate would have oppositve effects on them.
Therefore, the equilibrium with valued currency is not Pareto superior 
to the nonvalued- currency equilibrium.

\subsection{(d) Answer.}
After introducing government purchase ($L_g$) and ``inside money'' ($M_I$), 
the loan market clearing condition is given by 
\[
\frac{y}{2} N_1 - \frac{Y}{2 R_t} N_2 + L_g = \frac{H_0}{p_t} + \frac{M_I}{p_t}
\]
Since we assume that $L_g p_t = M_I$, the equation above is exactly the same as the one in (c).
Therefore, the equilibrium interest rate remains the same. 
(Namely, there is no inflation with inside money).

\noindent \textbf{Remark:} The inside money is not inflatory because assets are identical (in particular in terms of liquidity) 
and thus issuing the same amount of inside money as government purchse doesn't change the net supply of money.
See Kiyotaki-Moore (2018) for the discussion of Inside Money and Liquidity.

\section{Exercise 10.9}
\subsection{(a) - (b) Answers.}
\textbf{Definition}
The equilibrium with valued fiat money is given by prices ${p_t}$, and quantities 
$c^0_1, \{c_t^t \}, \{c_{t+1}^t \}$, such that 
\begin{enumerate}
    \item given by prices, the quanties solves all agents problem:
    that is, 
    for the initial old generation, $c^0_1$ solves 
    \begin{align*}
    \max_{c^0_1} \quad & \log u(c^0_1) \\ 
    \text{ subject to } \qquad & c^0_1 =  \frac{H}{p_1}
    \end{align*}
    For generation $t \geq 1$, $\{c^t_t\}, \{c^t_{t+1}\}, m_t$ solve 
    \begin{align*}
    \max_{c^t_t, c^t_{t+1}, m_t } \quad & \log u(c^t_t) +  \log u(c^t_{t+1}) \\ 
    \text{ subject to } \qquad & c^t_t + \frac{m_t}{p_t} =  y \\
    & c^t_{t+1}  =  \frac{m_t}{p_{t+1}}
    \end{align*}
    \item Market clears 
    \[
    \underbrace{N(y - c^t_t)}_{\text{saving from the young}} = \underbrace{\frac{H}{p_t}}_{\text{de-saving from the old}}.
    \]
\end{enumerate}
In equilibrium, the consumption is given by 
\[
c_t^t = \frac{y}{2},
\]
then the saving function is $s = \frac{y}{2}$,
therefore, using the market clearing condition, we have 
\[
N \frac{y}{2} = \frac{H}{p}
\]
As a result, the equilibrium price level is given by 
\[
p = \frac{2H}{Ny}.
\]

Now, let's introduce taxations. Then, the problem above becomes 
\begin{align*}
    \max_{c^t_t, c^t_{t+1}, m_t} \quad & \log u(c^t_t) +  \log u(c^t_{t+1}) \\ 
    \text{ subject to } \qquad & c^t_t + \frac{m_t}{p_t} =  y - \tau  \\
    & c^t_{t+1} = \frac{m_t}{p_{t+1}} +  \tau
\end{align*}
Then the Lagragian formulation is given by 
\[
\mathcal{L} = \log u(c^t_t) +  \log u(c^t_{t+1}) + \lambda \left[y - \tau - c_t^t - \frac{p_{t+1}}{p_t}(c_{t+1}^t - \tau )\right]
\]
Then the FOC implies $c^t_{t+1} = \frac{p_t}{p_{t+1}} c_t$.
Using the budget constraint, then we have 
\[
c_t^t + \frac{p_{t+1}}{p_t}(\frac{p_t}{p_{t+1}} c_t - \tau) = y - \tau
\]
It follows that 
\[
c_t^t = \frac{1 }{2}(y - \tau + \frac{p_{t+1}}{p_t} \tau) = \frac{1 }{2}(y - \tau + \frac{\tau}{R_t})
\]
Therefore, the aggregate saving is given by 
\[
N (y - \tau - c_t^t) = \frac{N}{2}(y - \tau - \frac{\tau}{R_t})
\]
Then using the market clearing condition and assuming a stationary price system, we have 
\[
\frac{N}{2}(y - 2 \tau) = \frac{H}{p} \Rightarrow p = \frac{2H}{N(y - 2 \tau)}.
\]

\subsection{(c) - (d) Answers.}
Notice that the equilibrium exists iff $\tau < \frac{y}{2}$.

It follows that $\frac{\partial p}{\partial \tau} > 0 $, namely, 
increasing tax burden will increase price level.


\section{Exercise 10.10}
\subsection{(a) Answer.}
Using the same setup as in Exercise 10.8 (a), we can derive the interest rate in the equilibrium without fiat money as
\[
1 + r_t = R_t = \frac{\beta N_2}{\alpha N_1}.
\]

\subsection{(b) Answer.}
Due to the lack of inter-generational trade, this equilibirum is not Pareto Optimal. (Yes?)


\subsection{(c) Answer.}
Using the saving function and the money market clearing condition, we can derive a first-order difference 
equation for the price level.
Concretely, it is given by 
\[
\frac{\alpha}{2} N_1 - \frac{\beta p_{t+1}}{2 p_t} N_2 = \frac{H_t}{p_t} \Rightarrow \alpha N_1 p_t - \beta N_2 p_{t+1} = 2 H_t
\]
Using Lag-operation, we have 
\[
\alpha N_1 p_t - \beta N_2 L^{-1} p_{t} = 2 H_t
\]
Then, we have 
\begin{align*}
p_t & = \left(1 - \frac{\beta N_2}{\alpha N_1} L^{-1}\right)^{-1} \frac{2H_t}{\alpha N_1} + \kappa \left(\frac{\alpha N_1}{\beta N_2}\right)^t \\ 
    & = \bigg[\frac{2H_0 }{\alpha N_1} \sum_{j = 0}^{\infty} \left(\frac{\beta N_2}{\alpha N_1}\right)^j L^{-j} z_t \bigg] + \kappa \left(\frac{\alpha N_1}{\beta N_2}\right)^t \\
    & = \bigg[\frac{2H_0 }{\alpha N_1} \sum_{j = 0}^{\infty} \left(\frac{\beta N_2}{\alpha N_1}\right)^j z_{t + j} \bigg] + \kappa \left(\frac{\alpha N_1}{\beta N_2}\right)^t \\
    & = \frac{2H_0 z^t }{\alpha N_1 - \beta N_2} + \kappa \left(\frac{\alpha N_1}{\beta N_2}\right)^t
\end{align*}
Therefore, there exists a continuum of equilibria parameterized by $\kappa$.

\subsection{(d) Answer.}
Let $\kappa = 0$, then we have ``quantity theory'' equation:
\[
p_t = \frac{2H_0 z^t }{\alpha N_1 - \beta N_2} = \frac{2H_t }{\alpha N_1 - \beta N_2}
\]


%% =========== %% ========= %%
\bibliography{../bib/references.bib}

\end{document}